\section{换一种界面读数}
\label{sec:interface}

第一套规则没有改。这一节是另一次计算：界面不再用前25\%高分区的二值标签，空分布也不再把坐标打乱成互不相关的点。细胞比例用四个标志基因程序的非负最小二乘，不是cell2location。Visium中心距约100\,$\mu$m，50\,$\mu$m和20--50\,$\mu$m的细胞邻接没有算。没有多重免疫荧光、共培养或类器官。

成纤维程序是DCN、LUM、COL1A1、FAP，巨噬程序是CD68、CSF1R、C1QA、C1QB，另有上皮和T细胞程序，用来把比例分开。被检验的配体和受体不进入这四个程序。界面权重是成纤维比例乘巨噬比例。点位信号仍是本点受体乘以150\,$\mu$m邻域的配体均值，切片读数改成权重下的加权平均。空分布把配体图平移200--500\,$\mu$m，200次，种子20261004。9例的汇总是中位数和有放回抽样2000次的区间，不是至少5/9例的投票。GC6-PM不进入这9例。

\begin{figure}[htbp]
  \centering
  \includegraphics[width=\textwidth]{fig_interface.pdf}
  \caption{平移空分布下的界面读数。A：每个点是一张原发灶切片，横轴是加权均值除以该切片空分布第95百分位，黑线是9例中位数，红虚线是1。B：蓝点是成纤维比例与150\,$\mu$m邻域巨噬比例的相关，红圈是同一张切片上把巨噬比例平移后的第95百分位。}
  \label{fig:interface}
\end{figure}

加权平均把旧分析里被中位数压成0的读数抬了起来，但没有抬到平移空分布之上（图~\ref{fig:interface}，表~\ref{tab:interface}）。C3--C3AR1的9例中位数是0.223，最小到最大是0.092--0.833；超额中位数是$-0.005$，抽样区间$-0.022$到$-0.002$，9例里8例低于自己的空分布。SPP1--CD44的超额在9/9例都为负，中位数$-0.018$，区间$-0.035$到$-0.003$。CCL2--CCR2的加权中位数只有0.012，97\%的界面权重落在信号为0的点位上，超额中位数$-0.001$。CD99--CD99的绝对读数有1.780，超额仍是$-0.021$，区间不跨过0。APP--CD74最高，中位数2.626，超额$-0.068$，区间$-0.136$到0.015，跨过0。SEMA4D--PLXNB2的区间也跨过0。TGFB1与两个受体分开算、再取TGFBR1和TGFBR2的较小值，三条读数的超额中位数都是负的。

\begin{table}[htbp]
  \centering
  \caption{9例原发灶、150\,$\mu$m、比例权重。超额是观察值减去平移空分布第95百分位。括号里是9例中位数经2000次有放回抽样的2.5\%和97.5\%分位。}
  \label{tab:interface}
  \footnotesize
  \setlength{\tabcolsep}{3pt}
  \begin{tabular}{lrrrr}
    \toprule
    配体受体 & 加权中位数 & 最小--最大 & 超额中位数 & 低于空分布 \\
    \midrule
    CCL2--CCR2 & 0.012 & 0.002--0.053 & $-0.001$ ($-0.002$--0.000) & 6/9 \\
    C3--C3AR1 & 0.223 & 0.092--0.833 & $-0.005$ ($-0.022$--$-0.002$) & 8/9 \\
    SPP1--CD44 & 0.088 & 0.007--0.758 & $-0.018$ ($-0.035$--$-0.003$) & 9/9 \\
    TGFB1--TGFBR1 & 0.100 & 0.026--0.198 & $-0.002$ ($-0.007$--0.001) & 5/9 \\
    TGFB1--TGFBR2 & 0.133 & 0.082--0.305 & $-0.006$ ($-0.012$--$-0.0002$) & 8/9 \\
    TGFB1--TGFBR复合物 & 0.060 & 0.014--0.078 & $-0.001$ ($-0.002$--0.001) & 6/9 \\
    APP--CD74 & 2.626 & 1.415--3.990 & $-0.068$ ($-0.136$--0.015) & 7/9 \\
    CD99--CD99 & 1.780 & 0.270--7.242 & $-0.021$ ($-0.032$--$-0.006$) & 8/9 \\
    SEMA4D--PLXNB2 & 0.277 & 0.174--0.723 & $-0.015$ ($-0.023$--0.002) & 7/9 \\
    \bottomrule
  \end{tabular}
\end{table}

半径改成100、200、300\,$\mu$m之后，超额中位数仍在0的负侧。唯一的例外是TGFB1--TGFBR1在100\,$\mu$m的超额中位数为$+0.0001$。把界面改回二值、取中位数，75\%分位以上同时富集成纤维和巨噬的点位上，C3--C3AR1、CCL2--CCR2和TGFB1--TGFBR1的9例中位数都是0。中位数会被零压住；加权平均避开了这个零，却仍不超过保留空间结构的空分布。

成纤维比例和邻域巨噬比例的相关，9例中位数是0.19。4/9例高于平移空分布，超额的中位数是$-0.022$。GC4的相关是0.36，空分布第95百分位是0.41，观察值仍低。第一套里GC4的C3--C3AR1中位数高于坐标打乱，在这一套平移空分布下没有成为稳定超额。

STAT3和SOCS3的转录均值，与C3--C3AR1、CCL2--CCR2、SPP1--CD44点位信号的Spearman中位数是0.08、0.02和0.15，没有一张切片超过0.21。这是转录本上的相关，不是磷酸化，也不是把一条边从网络里删掉之后的下游变化。

GC6-PM的C3--C3AR1加权均值是0.922，平移空分布第95百分位是0.925，超额$-0.004$。绝对水平和原发灶不是同一档，仍没有高于这张切片自己的空分布。它的转移部位记录不是腹膜，不并入9例。

单细胞按同一套程序归类。原发灶sample22只有6个成纤维细胞和18个上皮细胞，这两类不进入中位数，所以成纤维和上皮的原发分母是25份，巨噬和T细胞是26份。3份腹膜对原发只报方向（表~\ref{tab:source}）。

C3在原发灶里成纤维中位数最高，为0.374。3份腹膜肿瘤的成纤维中位数是2.032，3/3高于原发中位数；上皮是0.958对0.031，也是3/3。sample19和sample38里最高的一类是成纤维，sample20是上皮。SPP1在原发灶和腹膜里都以巨噬最高，腹膜巨噬中位数0.729，原发0.211，2/3高于原发。成纤维两边都低，0.068对0.053。sample19这一份里最高的一类是上皮。

CCL2在原发灶成纤维的中位数是1.084。3份腹膜成纤维中位数是0.422，0/3高于原发。前面样本均值里CCL2的腹膜方向，不是成纤维升高。腹膜里CCL2最高的类也不固定：sample19和sample38是上皮，sample20是巨噬。上皮中位数0.250对0.054，3/3高于原发，绝对值仍低于原发成纤维。TGFB1在3份腹膜里最高的类都是T细胞，中位数0.802对0.565，3/3；成纤维和巨噬都是1/3。

受体这边，C3AR1落在巨噬，腹膜中位数0.566，原发0.342，3/3。CCR2在腹膜T细胞的中位数是0.085，高于巨噬的0.045。CD44在巨噬和T细胞上都高，腹膜巨噬中位数1.076，低于原发巨噬1.251。

\begin{table}[htbp]
  \centering
  \caption{GSE183904按细胞类的样本均值。腹膜肿瘤3份，原发肿瘤成纤维和上皮25份、巨噬和T细胞26份。高于原发中位数的计数只作方向。}
  \label{tab:source}
  \footnotesize
  \setlength{\tabcolsep}{4pt}
  \begin{tabular}{llrrr}
    \toprule
    基因 & 细胞类 & 腹膜中位数 & 原发中位数 & 高于原发 \\
    \midrule
    C3 & 成纤维 & 2.032 & 0.374 & 3/3 \\
    C3 & 巨噬 & 0.834 & 0.085 & 2/3 \\
    C3 & 上皮 & 0.958 & 0.031 & 3/3 \\
    C3 & T细胞 & 0.087 & 0.003 & 3/3 \\
    SPP1 & 成纤维 & 0.068 & 0.053 & 2/3 \\
    SPP1 & 巨噬 & 0.729 & 0.211 & 2/3 \\
    SPP1 & 上皮 & 0.065 & 0.018 & 2/3 \\
    SPP1 & T细胞 & 0.020 & 0.009 & 3/3 \\
    CCL2 & 成纤维 & 0.422 & 1.084 & 0/3 \\
    CCL2 & 巨噬 & 0.217 & 0.290 & 1/3 \\
    CCL2 & 上皮 & 0.250 & 0.054 & 3/3 \\
    CCL2 & T细胞 & 0.021 & 0.008 & 3/3 \\
    TGFB1 & 成纤维 & 0.226 & 0.443 & 1/3 \\
    TGFB1 & 巨噬 & 0.480 & 0.535 & 1/3 \\
    TGFB1 & 上皮 & 0.327 & 0.230 & 2/3 \\
    TGFB1 & T细胞 & 0.802 & 0.565 & 3/3 \\
    \bottomrule
  \end{tabular}
\end{table}
